\documentclass[letterpaper,landscape,9pt,fleqn]{extarticle}
%\usepackage[utf8]{inputenc}
%\usepackage[T1]{fontenc}
\usepackage{graphicx}
\usepackage{xcolor}
\usepackage{tikz}
\usepackage{url}
\usepackage{qrcode}
%\usetikzlibrary{shapes.geometric}
%\usetikzlibrary{calc}
%\usepackage{fourier}
\usepackage{graphicx,nicefrac}
\usepackage{isomath,upgreek,xcolor,comment}
\usepackage{pdfpages}
%\usepackage{tkz-euclide}
%\usetkzobj{all}
\pagestyle{empty}
\usepackage[activate={true,nocompatibility},final,tracking=true,kerning=true,factor=1100,stretch=10,shrink=10]{microtype}
\usepackage[american]{babel}
\usepackage{centernot}
\usepackage{amstext} % for \text macro
\usepackage{array}   % for \newcolumntype macro
\newcolumntype{L}{>{$}l<{$}} % math-mode version of "l" column type

\newcommand{\dom}{\mathrm{dom}} 
\newcommand{\range}{\mathrm{range}} 
\newcommand{\codom}{\mathrm{codom}}
\newcommand{\zero}{\mathrm{zero}} 
\newcommand{\reals}{\mathbf{R}} 
\newcommand{\complex}{\mathbf{C}} 
\newcommand{\ball}{\mathrm{ball}}
\newcommand{\integers}{\mathbf{Z}} 
\newcommand{\im}{\mathrm{i}}
\newcommand{\euler}{\mathrm{e}}
\newcommand{\ssep}{\mid}
\newcommand{\arcsec}{\mathrm{arcsec}}
\newcommand{\arccsc}{\mathrm{arccsc}}
\newcommand{\arccot}{\mathrm{arccot}}
\newcommand{\glb}{\mathrm{glb}}
\newcommand{\lub}{\mathrm{lub}}
\newcommand{\length}{\mathrm{length}}
\usepackage{mathtools}
\DeclarePairedDelimiter{\parens}{\lparen}{\rparen}
\usepackage{amsmath,amssymb,textcomp}
\everymath{\displaystyle}

%\usepackage{times}
%\renewcommand\familydefault{\sfdefault}
%\usepackage{tgheros}
%\usepackage[defaultmono,scale=0.85]{droidmono}
%\usepackage{fourier}
\usepackage{multicol}
\setlength{\columnseprule}{0pt}
\setlength{\columnsep}{20.0pt}

\usepackage[]{enumerate}
\usepackage{expdlist}

\newenvironment{alphalist}{
  \begin{enumerate}[(a)]
    \addtolength{\itemsep}{-1.0\itemsep}}
  {\end{enumerate}}

\usepackage{geometry}
\geometry{letterpaper,left=0.4in,right=0.4in,top=0.4in,bottom=0.4in}

%\linespread{1.3}


% custom title
\makeatletter
\renewcommand*{\maketitle}{%
\noindent
\begin{minipage}{0.4\textwidth}
\begin{tikzpicture}
\node[rectangle,rounded corners=6pt,inner sep=10pt,fill=blue!50!black,text width= 0.95\textwidth] {\color{white}\Huge \@title};
\end{tikzpicture}
\end{minipage}
\hfill
\begin{minipage}{0.55\textwidth}
\begin{tikzpicture}
\node[rectangle,rounded corners=3pt,inner sep=10pt,draw=blue!50!black,text width= 0.95\textwidth] {\LARGE \@author};
\end{tikzpicture}
\end{minipage}
\bigskip\bigskip
}%
\makeatother

% custom section
\usepackage[explicit]{titlesec}
\newcommand*\sectionlabel{}
\titleformat{\section}
  {\gdef\sectionlabel{}
   \normalfont\sffamily\Large\bfseries\scshape}
  {\gdef\sectionlabel{\thesection\ }}{0pt}
  {
\noindent
\begin{tikzpicture}
\node[rectangle,rounded corners=3pt,inner sep=4pt,fill=blue!50!black,text width= 0.95\columnwidth] {\color{white}\sectionlabel#1};
\end{tikzpicture}
  }
\titlespacing*{\section}{0pt}{15pt}{10pt}


% custom footer
%\usepackage{fancyhdr}
%\makeatletter
%\pagestyle{fancy}
%\fancyhead{}
%\fancyfoot[C]{\footnotesize \textcopyright\ \@date\ \ \@author}
%\renewcommand{\headrulewidth}{0pt}
%\renewcommand{\footrulewidth}{0pt}
%\makeatother

\raggedbottom 
\raggedright
%\usepackage{tikz-3dplot}
\begin{document}

%\maketitle

\begin{multicols*}{3}

\section*{Greek characters}
\begin{tabular}{|L | L| L|} \hline
\mbox{Name} & \mbox{Symbol} & \mbox{Typical use(s)} \\ \hline
\mathrm{alpha} & \alpha  & \mbox{angle, constant} \\
\mathrm{beta} & \beta  & \mbox{angle, constant}  \\ 
\mathrm{gamma} & \gamma & \mbox{angle, constant} \\
\mathrm{delta} & \delta  & \mbox{limit definition}\\
\mathrm{epsilon} & \epsilon  \mbox{ or } \varepsilon & \mbox{limit definition} \\
\mathrm{theta}  & \theta  \mbox{ or } \vartheta &\mbox{angle}\\ 
\mathrm{pi} & \pi \mbox{ or } \uppi & \mbox{circular constant} \\
\mathrm{phi} & \phi \mbox{ or } \varphi  & \mbox{angle, constant} \\
\hline
\end{tabular}

\section*{Named Sets}
\begin{minipage}[l]{0.15\textwidth}
    \begin{tabular}{|L | L |} \hline 
        \mathrm{empty\,\, set} & \varnothing \\ 
        \mathrm{real\,\, numbers} & \mathbf{R} \\
        \mathrm{ordered \, \, pairs }  & \mathbf{R}^2 \\
        \hline
    \end{tabular}
\end{minipage}
    \begin{minipage}[l]{0.15\textwidth}
        \begin{tabular}{|L | L |} \hline 
            \mathrm{integers } & \mathbf{Z} \\
            \mathrm{positive \,\,  integers } & \mathbf{Z}_{>0} \\ 
            \mathrm{positive \,\, reals} & \mathbf{R}_{>0} \\
            \hline
        \end{tabular}              
\end{minipage}

\section*{Set Symbols}
\begin{minipage}[l]{0.15\textwidth}
    \begin{tabular}{|L | L|} \hline
        \mbox{Meaning}  & \mbox{Symbol} \\ \hline
        \mathrm{is \,\, a \,\, member} & \in \\
        \mathrm{subset}       & \subset \\
        \mathrm{intersection} & \cap \\ \hline
        \end{tabular}   
\end{minipage}
\begin{minipage}[l]{0.15\textwidth}
    \begin{tabular}{|L | L|} \hline
        \mbox{Meaning}  & \mbox{Symbol} \\ \hline
        \mathrm{union} & \cup  \\ 
        \mathrm{complement} & \mbox{superscript}^\mathrm{C} \\
        \mathrm{set \,\,  minus}  & \setminus \\ \hline
        \end{tabular}   
\end{minipage}

\section*{Logic Symbols}
\begin{minipage}[l]{0.15\textwidth}
    \begin{tabular}{|L | L|} \hline 
        \mbox{Meaning}  & \mbox{Symbol} \\ \hline 
        \mathrm{negation} &  \lnot   \\
        \mathrm{and} &  \land  \\
        \mathrm{or} &  \lor  \\
        \mathrm{implies} &  \implies \\ \hline    
    \end{tabular}   
\end{minipage}
\begin{minipage}[c]{0.15\textwidth}
    \begin{tabular}{|L | L|} \hline 
        \mbox{Meaning}  & \mbox{Symbol} \\ \hline 
        \mbox{equivalent} &  \equiv \\ 
        \mbox{iff} & \iff \\ 
        \mbox{for all} & \forall \\
        \mbox{there exists} & \exists \\ \hline
    \end{tabular}   
\end{minipage}

\section*{Function Notation}
\begin{tabular}{|L | L|} \hline 
    \dom(F) &   \mbox{domain of function } F \\
    \range(F) &   \mbox{range of function } F \\
    \mathrm{C}_{A} & \mbox{ set of continuous functions on set } A \\
    \mathrm{C}_{A}^1 & \mbox{ set of differentiable functions on set } A \\
    A \to B   & \mbox{set of functions from } A \mbox { to } B \\ \hline
\end{tabular}

\section*{Magnitude \& Conjugate}
For all \(a,b \in \reals\)
\begin{alignat*}{1}
   |a + \im b| &= \sqrt{a^2+b^2} \\
   \overline{a + \im b} &= a - \im b
\end{alignat*}
For all \(x,y,z \in \complex\), we have
\begin{alignat*}{1}
   &|x y| = |x| |y| \\
    &|x+y| \leq |x| + |y| \\
    &\big | |x| - |y|  \big |  \leq |x-y|    \\
    &\overline{xy} = \overline{x} \overline{y} \\
    &\frac{1}{z} = \frac{\overline{z}}{|z|^2} (\mbox{ for } z \neq 0) \\
    &\overline{\left(\frac{x}{y} \right)} = \frac{\overline{x}} {\overline{y}} \\
\end{alignat*}

\section*{Complex Exponential}
For $x,y \in \reals$
\begin{alignat*}{1}
  &\euler^{\im y} = \cos(y) + \im \sin(y) \\
  &\euler^{x + \im y} = \euler^{x} \left(\cos(y) + \im \sin(y) \right)
\end{alignat*}
For all $z_1, z_2 \in \complex$,
\begin{alignat*}{1}
    &\euler^{z_1 + z_2} = \euler^{z_1} \euler^{z_2}\\
    &\left[\euler^{z_1} =  \euler^{z_2} \right]
     \equiv \left[z_1 - z_2 = 2 \uppi n, n \in \integers \right]
\end{alignat*}

\section*{Argument \& Polar form}
For all \(z \in \complex_{\neq 0}\), there is a unique \(\theta \in (-\uppi, \uppi]\) such that
\begin{alignat*}{1}
   &z = |z| (\cos(\theta) + \im \sin(\theta)) \\
   &\arg(z) = \theta
\end{alignat*}
We have
\begin{alignat*}{1}
  \sqrt{z} &= \sqrt{|z|} (\cos(\theta/2) + \im \sin(\theta/2)) \\
  z^a &= |z|^a  (\cos(\theta/a) + \im \sin(\theta/a))  (\mbox{ for } z \in \reals_{\neq 0}) \\
  \log(z) &= \log(|z|) + \im \arg(z)
\end{alignat*}

\section*{Topology \& Disks}
For $a \in \complex$ and $r \in \reals_{>0}$, we define
\begin{align*}
  D[a,r] &= \{z \in \complex | |z-a| < r \}, \\
  D^\prime[a,r] &= \{z \in \complex | 0 < |z-a| < r \}, \\
  C[a,r] &= \{z \in \complex | |z-a| = r \}.
\end{align*}

Let $G$ be a subset of $\complex$. We say 
\begin{itemize}
  \item $z$ is an \emph{interior point of} $G$ provided
  \((\exists r \in \reals_{>0})(D[z,r] \subset G) \).

  \item $z$ is a \emph{boundary point of} $G$ provided
  \((\forall r \in \reals_{>0})(D[z,r] \cap G \neq \varnothing \land
  D[z,r] \cap G^C \neq \varnothing \}\).

  \item $z$ is an \emph{acculumation point of} $G$ provided
  \((\forall r \in \reals_{>0})(D^\prime [z,r] \cap G \neq \varnothing \}\).

  
\end{itemize}

\section*{Limits}

Let $G$ be a subset of $\complex$, $z_o$ be an accumulation point of $G$, and $F \in G \to \complex$.
We say $F$ \emph{has a limit towards} $z_o$ provided
\begin{alphalist}
\item $\exists L \in \complex$
\item $\forall \varepsilon \in \reals_{>0}$
\item $\exists \delta  \in \reals_{>0}$
\item $\forall z \in D[z_o, \delta] $
\item we have $|F(z) - L| < \varepsilon$.
\end{alphalist}

\section*{Continuity}

Let $G$ be a subset of $\complex$ and let $z_o \in \dom{F}$. We say $F$ \emph{is continuous at} $z_o$ provided
\begin{alphalist}
\item $\forall \varepsilon \in \reals_{>0}$
\item $\exists \delta  \in \reals_{>0}$
\item $\forall z \in D[z_o, \delta] $
\item we have $|F(z) - F(z_o)| < \varepsilon$.
\end{alphalist}

\section*{Derivatives}
Let $G$ be a subset of $\complex$ and let $z_o$ be an interior point of $G$. We say $F$ \emph{is differentiable at} $z_o$ provided 
\begin{equation*}
  \lim_{z\to z_o} \frac{F(z) - F(z_o)}{z-z_o}
\end{equation*}
exists.
\vfill
\noindent {Revised \today. Barton Willis is the author of this work. This work is
licensed under Attribution 4.0 International (CC BY 4.0) \,  \qrcode[height=0.15in]{https://creativecommons.org/licenses/by/4.0/}. For the current version of
this document, visit \, \qrcode[height=0.15in]{https://github.com/barton-willis/Math-100-200-level}}
\end{multicols*}%{3}

\end{document}